% Generated by report/make_restated.py from theory/A1_lower_bound.tex: do not edit (edit the note and rerun).
\begin{restated}{Remark}{rem:plot}{A1}{0}{what to expect in a plot}
For $\varepsilon=10^{-6}$ ($\log(1/\varepsilon)\approx13.8$), $n=100$ and $m=200$: the dimension condition gives $1+n(\frac{16m^2}{\log^2(1/\varepsilon)}-1)\approx3.4\cdot10^5$ and the slow-mode condition $1+\frac{n(m+1)^2}{\pi^2}\approx4.1\cdot10^5$, so $\kappa_{\max}\approx3.4\cdot10^5$ and the window $401\le\kappa\le3.4\cdot10^5$ of Proposition~\ref{prop:A4}(ii) spans almost three decades, enough to read the slopes $1$ and $\frac12$, with the plateau below $\kappa\approx n$ and the bend above $\kappa_{\max}$. Two caveats. First, a sandwich $c_1\kappa\le T\le c_2\kappa$ pins the slope of a line fitted over $D$ decades only up to $\pm\log_{10}(c_2/c_1)/D$; for SAGA, $c_2/c_1=\frac{16}3\cdot\frac{\log((1+n/4)/\varepsilon)}{\log(\rho/\varepsilon)}\approx7.7$ with the smallest value of the exact $\rho$ over the window, $\rho\approx0.13$ (reached just above $\kappa\approx2.2\cdot10^5$, Remark~\ref{rem:rho}), and for L-Katyusha the ratio of the two sides of \eqref{eq:lkwindow} is between $91$ (bottom of the window, $C_0$ bounded by $1+\frac58Q$ at that point) and about $128$ ($C_0$ replaced by its bound $1+\frac58Q$ at the top of the window, which gives constants valid over the whole window). The order of growth is a theorem; the numerical value of the slope must be read from exact quantities: the mean-iterate time $t^\ast(\kappa)=\min\{t:\norm{(I-\gamma\nabla^2F)^t(\theta_0-\etastar)}^2\le\varepsilon\norm{\theta_0-\etastar}^2\}$, computable from the $m$ eigenpairs of $A_m(\zeta_Q)$, is a lower bound on $T_{\rm SAGA}$ (first inequality of \eqref{ineq:saga-mean}), and for these parameters we found $t^\ast\approx1.4\,\kappa\log(1/\varepsilon)$ with log-log slope $0.99$ over $\kappa\in[10^3,3\cdot10^5]$. Second, with $m=20$ the window collapses ($\kappa_{\max}\approx3.3\cdot10^3$, less than a decade above $4n+1$): the slopes measured over $\kappa\in[10^3,10^6]$ would be far below $1$ and $\frac12$. This finite-dimension effect is quantified by \eqref{eq:kappaeff}: plotting $T_{\rm SAGA}$ against $\kappa_{\rm eff}$ instead of $\kappa$ restores a linear growth (Corollary~\ref{cor:saga} with $\nu'=\mu_F$), which is the cleanest test of the ``effective constants'' idea of hypothesis (B2) on an instance where the infimum over $\theta$ of the effective curvature $\muEff(\theta)=\inner{\theta-\etastar}{\Fp(\theta)}/\norm{\theta-\etastar}^2$ of (B2) is known exactly: it is $\mu_F=\lmin(\nabla^2F)$. For the configuration of the experiments of the report ($n=16$, $m=256$, $\varepsilon=10^{-8}$), the dimension condition gives $\kappa_{\max}\approx4.9\cdot10^4$ (the slow-mode condition gives $1.1\cdot10^5$), so the window $65\le\kappa\le4.9\cdot10^4$ spans $2.9$ decades, the smallest exact $\rho$ in it is $0.15$ and $c_2/c_1\approx6.5$ for SAGA; the swept values $\kappa\ge10^5$ lie outside the window, and slopes should be fitted over $10n\le\kappa\le\kappa_{\max}$.
\end{restated}
